% ======================================================================
\chapter[Постулаты Бора. Атом водорода. Водородоподобные атомы]%
        {Постулаты Бора. Атом водорода и водородоподобные системы}
\label{t:03}
% ======================================================================

\begin{qbox}
\textbf{Билет 3.} Постулаты Бора для атома водорода. Характерные параметры для атома
водорода: боровский радиус, скорость электрона, уровни энергии. Водородоподобные атомы:
влияние массы ядра (изотопический сдвиг), атомы щелочных металлов, многозарядные ионы, мезоатомы.
\end{qbox}

\section{Содержание билета}

\textbf{Подтемы:}
\begin{enumerate}
 \item постулаты Бора; правило квантования момента импульса;
 \item радиусы орбит, скорости, энергии уровней; боровский радиус, постоянная тонкой структуры;
 \item спектральные серии водорода, формула Ридберга;
 \item учёт конечной массы ядра: приведённая масса, изотопический сдвиг (H и D);
 \item атомы щелочных металлов: квантовый дефект;
 \item многозарядные (водородоподобные) ионы: масштабирование по $Z$;
 \item мезоатомы (мюонные атомы), позитроний.
\end{enumerate}

\textbf{Необходимые понятия:} кулоновское взаимодействие, движение по окружности,
задача двух тел и центр масс, фотон $\hbar\omega$ (\tref{t:01}), волна де Бройля (\tref{t:02}).

\section{Определения и физические принципы}

\begin{principle}{Постулаты Бора (1913)}
\begin{enumerate}[leftmargin=1.4em]
 \item \emph{Постулат стационарных состояний.} Атом может длительно находиться только в
       определённых стационарных состояниях с дискретными энергиями $E_1<E_2<\dots$; находясь
       в них, атом не излучает, несмотря на ускоренное движение электрона.
 \item \emph{Правило частот.} При переходе атома из состояния $n$ в состояние $m$ испускается
       (или поглощается) фотон с частотой
       \begin{equation}\label{eq:3:freq}
          \hbar\omega_{nm}=E_n-E_m .
       \end{equation}
 \item \emph{Правило квантования} (для круговых орбит). В стационарных состояниях момент
       импульса электрона кратен $\hbar$:
       \begin{equation}\label{eq:3:quant}
          L=\mu vr=n\hbar,\qquad n=1,2,3,\dots
       \end{equation}
\end{enumerate}
\end{principle}

\emph{Физический смысл и проблема.} По классической электродинамике электрон, движущийся с
ускорением по орбите, должен излучать и за время $\sim10^{-11}$~с упасть на ядро; спектр
излучения должен быть сплошным. Опыт показывает устойчивость атомов и линейчатые спектры.
Постулаты Бора \emph{вводят} дискретность как новый принцип, не объясняя её. Правило
квантования~\eqref{eq:3:quant} эквивалентно условию, что на орбите укладывается целое число
волн де Бройля: $2\pi r=n\lambda=nh/(\mu v)$.

\begin{definition}[характерные атомные величины]
\begin{align*}
 &\text{боровский радиус}&& \aB=\frac{\hbar^2}{\mel e^2}=0{,}529\ \text{Å},\\
 &\text{постоянная тонкой структуры}&& \alpha=\frac{e^2}{\hbar c}\approx\frac{1}{137},\\
 &\text{энергия Ридберга}&& \Ry=\frac{\mel e^4}{2\hbar^2}=\frac{e^2}{2\aB}=\frac{\alpha^2\mel c^2}{2}=13{,}6\eV .
\end{align*}
\end{definition}

\begin{definition}[водородоподобная система]
Система из одного электрона (или другой отрицательно заряженной частицы) и ядра с зарядом $Ze$,
либо система, в которой один внешний электрон движется в поле, близком к кулоновскому
(атомы щелочных металлов).
\end{definition}

\emph{Область применимости модели Бора.} Модель даёт точные уровни энергии для кулоновской
задачи (совпадающие с решением уравнения Шредингера, \tref{t:10}), если движение
нерелятивистское ($Z\alpha\ll1$). Модель не объясняет интенсивности линий, не применима к
многоэлектронным атомам (уже к He), неверно предсказывает момент импульса основного
состояния ($\hbar$ вместо 0) и противоречит соотношению неопределённостей (точные орбиты).

\section{Основная теория}

\subsection{Уровни энергии атома водорода}

Решение системы «второй закон Ньютона + правило квантования» (вывод~\ref{der:3:bohr})
для частицы с зарядом $-e$ и приведённой массой $\mu$ в поле ядра с зарядом $Ze$:
\begin{equation}\label{eq:3:levels}
  r_n=\frac{\hbar^2}{\mu Ze^2}\,n^2,\qquad
  v_n=\frac{Ze^2}{n\hbar}=\frac{Z\alpha c}{n},\qquad
  E_n=-\frac{\mu Z^2e^4}{2\hbar^2}\frac{1}{n^2}.
\end{equation}
Для водорода ($Z=1$, $\mu\approx\mel$):
\begin{center}
\begin{tabular}{@{}ll@{}}
\toprule
радиус первой орбиты & $r_1=\aB=0{,}529$~Å \\
скорость на первой орбите & $v_1=\alpha c=2{,}19\cdot10^8$~см/с $\approx c/137$ \\
энергия основного состояния & $E_1=-\Ry=-13{,}6$~эВ (энергия ионизации 13{,}6~эВ) \\
период обращения & $T_1=2\pi\aB/v_1\approx1{,}5\cdot10^{-16}$~с \\
напряжённость поля ядра на орбите & $e/\aB^2\approx5\cdot10^9$~В/см \\
\bottomrule
\end{tabular}
\end{center}
Малость $v_1/c=\alpha$ оправдывает нерелятивистское описание; релятивистские поправки к уровням
имеют относительную величину $\sim\alpha^2\sim5\cdot10^{-5}$ (тонкая структура).

\subsection{Спектральные серии}

Из~\eqref{eq:3:freq} и~\eqref{eq:3:levels} (вывод~\ref{der:3:rydberg}):
\begin{equation}\label{eq:3:rydberg}
  \frac{1}{\lambda_{nm}}=R\,Z^2\left(\frac{1}{m^2}-\frac{1}{n^2}\right),\qquad n>m,
\end{equation}
где $R=\mu e^4/(4\pi\hbar^3c)$ — постоянная Ридберга. Линии с общим нижним уровнем $m$
образуют серию (рис.~\ref{fig:3:levels}): $m=1$ — серия Лаймана (УФ, 91--122~нм),
$m=2$ — серия Бальмера (видимая, $\mathrm H_\alpha=656{,}3$~нм, граница 364,6~нм),
$m=3$ — Пашена (ИК), $m=4$ — Брэккета, $m=5$ — Пфунда. При $n\to\infty$ линии сгущаются к
\emph{границе серии}, за которой лежит сплошной спектр (переходы с ионизацией).

\begin{figure}[htbp]
\centering
\begin{tikzpicture}[yscale=0.36, xscale=1.0]
  \foreach \n in {1,2,3,4,5} { \draw[thick] (0,{-13.6/(\n*\n)}) -- (10,{-13.6/(\n*\n)}); }
  \foreach \n/\val in {1/{-13,6},2/{-3,40},3/{-1,51}} {
     \node[left, font=\small] at (0,{-13.6/(\n*\n)}) {$n=\n$};
     \node[right, font=\scriptsize] at (10,{-13.6/(\n*\n)}) {$\val$ эВ};
  }
  \node[left, font=\small] at (0,-0.75) {$n=4,5,\dots$};
  \draw[thick, dashed] (0,0) -- (10,0) node[right, font=\scriptsize] {$0$};
  \node[left, font=\small] at (0,0.75) {$n=\infty$};
  \foreach \n in {2,3,4,5} { \draw[accent, ->] ({0.25+0.38*\n},{-13.6/(\n*\n)}) -- ({0.25+0.38*\n},-13.6); }
  \node[accent, font=\small, anchor=east] at (0.95,-8.0) {Лайман};
  \foreach \n in {3,4,5} { \draw[keycol, ->] ({3.7+0.42*\n},{-13.6/(\n*\n)}) -- ({3.7+0.42*\n},-3.4); }
  \node[keycol, font=\small] at (5.4,-4.6) {Бальмер};
  \foreach \n in {4,5} { \draw[dercol, ->] ({6.6+0.45*\n},{-13.6/(\n*\n)}) -- ({6.6+0.45*\n},{-13.6/9}); }
  \node[dercol, font=\small] at (8.6,-2.3) {Пашен};
\end{tikzpicture}
\caption{Уровни энергии атома водорода $E_n=-13{,}6/n^2$~эВ (в масштабе) и переходы,
образующие серии Лаймана, Бальмера и Пашена.}
\label{fig:3:levels}
\end{figure}

\subsection{Влияние массы ядра. Изотопический сдвиг}

Ядро имеет конечную массу $M$, поэтому электрон и ядро вращаются вокруг общего центра масс.
Задача двух тел сводится к движению одной частицы с \emph{приведённой массой}
\begin{equation}
  \mu=\frac{\mel M}{\mel+M}=\frac{\mel}{1+\mel/M}
\end{equation}
по относительной координате (вывод~\ref{der:3:reduced}). Все формулы~\eqref{eq:3:levels}
сохраняются с заменой $\mel\to\mu$; в частности,
\begin{equation}\label{eq:3:RM}
  R_M=\frac{R_\infty}{1+\mel/M},\qquad R_\infty=\frac{\mel e^4}{4\pi\hbar^3c}=109\,737\ \text{см}^{-1}.
\end{equation}
Для водорода $R_{\mathrm H}=109\,678$~см$^{-1}$. Изотопы одного элемента имеют разные $M$, и их
линии сдвинуты: \emph{изотопический сдвиг}. Для водорода и дейтерия ($M_D\approx2M_p$)
(вывод~\ref{der:3:isotope})
\begin{equation}\label{eq:3:shift}
  \frac{\lambda_{\mathrm H}-\lambda_{\mathrm D}}{\lambda_{\mathrm H}}
  \approx\frac{\mel}{M_p}-\frac{\mel}{M_D}\approx\frac{\mel}{2M_p}=2{,}7\cdot10^{-4}.
\end{equation}
Для линии $\mathrm H_\alpha$ ($656{,}3$~нм) сдвиг $\approx0{,}18$~нм — линия дейтерия смещена в
коротковолновую сторону. По такому слабому спутнику линий Бальмера Г.~Юри в 1932~г. открыл
дейтерий. (У тяжёлых атомов изотопический сдвиг определяется в основном не массой, а различием
размеров ядер — «объёмный эффект».)

\subsection{Атомы щелочных металлов}

В атомах Li, Na, K, Rb, Cs один валентный электрон находится вне заполненных оболочек
(«атомного остатка» с зарядом $+e$). Вдали от остатка электрон движется в кулоновском поле
заряда $+e$, но вблизи — в поле не полностью экранированного ядра. Отклонение от кулоновского
поля тем сильнее, чем глубже электрон проникает в остаток, т.\,е. чем меньше его орбитальный
момент $l$ (\tref{t:09}). Эмпирически уровни описываются формулой Ридберга
\begin{equation}\label{eq:3:qd}
  E_{nl}=-\frac{\Ry}{(n-\Delta_l)^2}=-\frac{\Ry}{n^{*2}},
\end{equation}
где $\Delta_l$ — \emph{квантовый дефект}, почти не зависящий от $n$, но зависящий от $l$; $n^*$ —
эффективное главное квантовое число. Для натрия $\Delta_s\approx1{,}37$, $\Delta_p\approx0{,}88$,
$\Delta_d\approx0{,}01$. Отсюда $E_{3s}\approx-13{,}6/(1{,}63)^2\approx-5{,}1$~эВ (потенциал ионизации
Na равен 5,14~эВ), а переход $3p\to3s$ даёт $\approx2{,}1$~эВ — жёлтая D-линия
($\lambda\approx589$~нм). Вырождение по $l$, характерное для чистого кулоновского поля, снимается.
Причины $\Delta_l\neq0$: проникновение электрона в остов и поляризация остова.

\subsection{Многозарядные ионы}

Ионы He$^+$, Li$^{2+}$, Be$^{3+}$,\,\dots, U$^{91+}$ имеют один электрон и ядро с зарядом $Ze$.
По~\eqref{eq:3:levels}
\begin{equation}
  r_n=\frac{\aB n^2}{Z},\qquad v_n=\frac{Z\alpha c}{n},\qquad E_n=-\frac{Z^2\Ry}{n^2}.
\end{equation}
Размеры уменьшаются в $Z$ раз, энергии растут в $Z^2$ раз. Для He$^+$ $E_1=-54{,}4$~эВ; серия
Пикеринга иона He$^+$ (переходы на $n=4$) содержит линии, почти совпадающие с линиями Бальмера,
так как $Z^2/n^2=4/(2k)^2=1/k^2$; небольшое различие обусловлено разными приведёнными массами.
С ростом $Z$ скорость $v_1=Z\alpha c$ становится сравнимой с $c$ (для U$^{91+}$ $v_1\approx0{,}67c$,
$|E_1|\approx115$~кэВ), и необходим релятивистский расчёт. Сходный закон $Z^2$ действует для
рентгеновских K-линий тяжёлых атомов: электрон, переходящий на вакансию в K-оболочке, «видит»
заряд $(Z-1)e$, что даёт закон Мозли $\hbar\omega_{K_\alpha}\approx\frac34\Ry\,(Z-1)^2$.

\subsection{Мезоатомы и экзотические атомы}

Отрицательный мюон ($m_\mu=206{,}8\,\mel$, время жизни 2,2~мкс), захваченный ядром, образует
\emph{мюонный атом}. Мюон каскадно переходит на нижние уровни за время $\sim10^{-13}$~с, много
меньшее времени жизни. В~\eqref{eq:3:levels} масса $\mel$ заменяется приведённой массой
$\mu=m_\mu M/(m_\mu+M)$:
\begin{itemize}
 \item мюонный водород ($\mu p$): $\mu=185{,}8\,\mel$, $r_1=\aB/185{,}8\approx285$~фм,
       $E_1=-185{,}8\cdot13{,}6\eV\approx-2{,}5$~кэВ;
 \item тяжёлые ядра: для свинца ($Z=82$) $r_1\approx\aB\mel/(207\cdot82)\approx3$~фм — меньше
       радиуса ядра ($\approx7$~фм). Мюон значительную часть времени проводит внутри ядра, и
       его уровни чувствительны к распределению заряда в ядре. Измерение рентгеновских линий
       мюонных атомов — один из методов определения размеров ядер (\tref{t:17}).
\end{itemize}
Существуют также пионные, каонные, антипротонные атомы. \emph{Позитроний} ($e^+e^-$) —
водородоподобная система с $\mu=\mel/2$: $E_1=-6{,}8$~эВ, расстояние между частицами в
основном состоянии вдвое больше $\aB$.

\section{Математические выводы}

\begin{derivation}[label=der:3:bohr]{радиусы, скорости и энергии стационарных орбит}
\emph{Допущения:} ядро с зарядом $Ze$; движение нерелятивистское, по окружности; масса
частицы $\mu$ (приведённая, см. вывод~\ref{der:3:reduced}).

\emph{Исходные уравнения:} второй закон Ньютона (кулоновская сила — центростремительная)
и правило квантования:
\begin{equation}\label{eq:3:sys}
  \frac{\mu v^2}{r}=\frac{Ze^2}{r^2},\qquad \mu vr=n\hbar .
\end{equation}
\emph{Шаг 1.} Умножим первое уравнение на $r^2$: $\mu v^2r=Ze^2$. Разделим на второе:
\[
  \frac{\mu v^2 r}{\mu vr}=\frac{Ze^2}{n\hbar}\quad\Longrightarrow\quad v_n=\frac{Ze^2}{n\hbar}
  =\frac{e^2}{\hbar c}\frac{Zc}{n}=\frac{Z\alpha c}{n}.
\]
\emph{Шаг 2.} Из правила квантования
\[
  r_n=\frac{n\hbar}{\mu v_n}=\frac{n\hbar\cdot n\hbar}{\mu Ze^2}=\frac{\hbar^2}{\mu Ze^2}n^2 .
\]
\emph{Шаг 3.} Полная энергия $E=T+U$, $T=\mu v^2/2$, $U=-Ze^2/r$. Из первого уравнения~\eqref{eq:3:sys}
$\mu v^2=Ze^2/r$, поэтому
\[
  T=\frac{Ze^2}{2r}=-\frac{U}{2},\qquad E=T+U=-\frac{Ze^2}{2r}=-T .
\]
(Это частный случай теоремы вириала для кулоновского поля: $2\bar T=-\bar U$.)
\emph{Шаг 4.} Подставим $r_n$:
\[
  E_n=-\frac{Ze^2}{2}\cdot\frac{\mu Ze^2}{\hbar^2n^2}=-\frac{\mu Z^2e^4}{2\hbar^2}\frac{1}{n^2}.
\]
\emph{Проверка размерности:} $[\mu e^4/\hbar^2]=\text{г}\cdot(\text{эрг}\cdot\text{см})^2/(\text{эрг}\cdot\text{с})^2=\text{г}\,\text{см}^2/\text{с}^2=\text{эрг}$;
$[\hbar^2/\mu e^2]=(\text{эрг}\,\text{с})^2/(\text{г}\cdot\text{эрг}\,\text{см})=\text{см}$.
\emph{Численно} ($Z=1$, $\mu=\mel$):
\[
 \aB=\frac{(\hbar c)^2}{\mel c^2\,e^2}=\frac{(1973{,}3\ \text{эВ}\cdot\text{Å})^2}{0{,}511\cdot10^6\ \text{эВ}\cdot14{,}40\ \text{эВ}\cdot\text{Å}}=0{,}529\ \text{Å},
\]
\[
 \Ry=\frac{e^2}{2\aB}=\frac{14{,}40}{2\cdot0{,}529}\ \text{эВ}=13{,}6\ \text{эВ}.
\]
\end{derivation}

\begin{derivation}[label=der:3:rydberg]{формула Ридберга}
По правилу частот~\eqref{eq:3:freq} и формуле~\eqref{eq:3:levels}:
\[
  \hbar\omega_{nm}=E_n-E_m=\frac{\mu Z^2e^4}{2\hbar^2}\Bigl(\frac{1}{m^2}-\frac{1}{n^2}\Bigr).
\]
С учётом $\omega=2\pi c/\lambda$:
\[
  \frac{1}{\lambda}=\frac{\omega}{2\pi c}=\frac{\mu e^4}{4\pi\hbar^3c}\,Z^2\Bigl(\frac{1}{m^2}-\frac{1}{n^2}\Bigr)
  \equiv R\,Z^2\Bigl(\frac{1}{m^2}-\frac{1}{n^2}\Bigr).
\]
Пример: $\mathrm H_\alpha$ ($3\to2$): $1/\lambda=R_{\mathrm H}(1/4-1/9)=R_{\mathrm H}\cdot5/36$,
$\lambda=36/(5\cdot109\,678\ \text{см}^{-1})=6{,}565\cdot10^{-5}$~см (в вакууме; в воздухе 656,3~нм).
\end{derivation}

\begin{derivation}[label=der:3:reduced]{приведённая масса}
Пусть $\vect r_e$, $\vect r_N$ — радиус-векторы электрона и ядра, потенциал зависит только от
$\vect r=\vect r_e-\vect r_N$. Введём центр масс $\vect R=(\mel\vect r_e+M\vect r_N)/(\mel+M)$.
Тогда $\vect r_e=\vect R+\frac{M}{\mel+M}\vect r$, $\vect r_N=\vect R-\frac{\mel}{\mel+M}\vect r$, и
кинетическая энергия
\[
  T=\frac{\mel\dot{\vect r}_e^{\,2}}{2}+\frac{M\dot{\vect r}_N^{\,2}}{2}
   =\frac{(\mel+M)\dot{\vect R}^2}{2}+\frac12\Bigl(\frac{\mel M^2}{(\mel+M)^2}+\frac{M\mel^2}{(\mel+M)^2}\Bigr)\dot{\vect r}^{\,2}
   =\frac{(\mel+M)\dot{\vect R}^2}{2}+\frac{\mu\dot{\vect r}^{\,2}}{2}
\]
(перекрёстные члены $\propto\dot{\vect R}\cdot\dot{\vect r}$ сокращаются). В системе центра масс
($\dot{\vect R}=0$) задача сводится к движению одной частицы массы $\mu=\mel M/(\mel+M)$ в поле
$U(r)$; полный момент импульса равен $\mu r v$, где $v=|\dot{\vect r}|$. Поэтому в~\eqref{eq:3:sys}
стоит $\mu$, а уровни $E_n\propto\mu$.
\end{derivation}

\begin{derivation}[label=der:3:isotope]{изотопический сдвиг H--D}
Для линии одного и того же перехода $\lambda\propto1/R_M\propto1+\mel/M$. Тогда
\begin{align*}
  \frac{\lambda_{\mathrm D}}{\lambda_{\mathrm H}}&=\frac{1+\mel/M_D}{1+\mel/M_p}
  \approx1-\frac{\mel}{M_p}+\frac{\mel}{M_D},\\
  \frac{\lambda_{\mathrm H}-\lambda_{\mathrm D}}{\lambda_{\mathrm H}}&\approx\frac{\mel}{M_p}-\frac{\mel}{M_D}
  \approx\frac{\mel}{M_p}\Bigl(1-\frac12\Bigr)=\frac{1}{2\cdot1836}=2{,}72\cdot10^{-4},
\end{align*}
где отброшены члены порядка $(\mel/M)^2$. Для $\mathrm H_\alpha$: $\Delta\lambda=656{,}3\cdot2{,}72\cdot10^{-4}\approx0{,}18$~нм.
\end{derivation}

\begin{derivation}[label=der:3:corr]{принцип соответствия для больших $n$}
Частота излучения при переходе $n\to n-1$ при $n\gg1$:
\[
  \omega_{n,n-1}=\frac{E_n-E_{n-1}}{\hbar}\approx\frac{1}{\hbar}\frac{\dd E_n}{\dd n}
  =\frac{1}{\hbar}\cdot\frac{\mu Z^2e^4}{\hbar^2n^3}=\frac{\mu Z^2e^4}{\hbar^3n^3}.
\]
Классическая частота обращения по $n$-й орбите:
$\omega_{\text{орб}}=\dfrac{v_n}{r_n}=\dfrac{Ze^2}{n\hbar}\cdot\dfrac{\mu Ze^2}{\hbar^2n^2}=\dfrac{\mu Z^2e^4}{\hbar^3n^3}$.
Частоты совпадают: при больших квантовых числах квантовая теория переходит в классическую
(принцип соответствия Бора).
\end{derivation}

\section{Итоговые формулы}

\keyformula{r_n=\frac{\aB}{Z}\frac{\mel}{\mu}n^2,\qquad v_n=\frac{Z\alpha c}{n},\qquad E_n=-\frac{Z^2}{n^2}\frac{\mu}{\mel}\Ry}
 {$\aB=\hbar^2/\mel e^2=0{,}529$~Å, $\alpha=e^2/\hbar c$, $\Ry=\mel e^4/2\hbar^2=13{,}6$~эВ, $\mu$ — приведённая масса, $n=1,2,\dots$}
 {одна частица в кулоновском поле ядра; $Z\alpha\ll1$.}
 {размеры, скорости и уровни водородоподобных систем (атомы, ионы, мезоатомы, позитроний).}

\keyformula{\frac1\lambda=RZ^2\Bigl(\frac1{m^2}-\frac1{n^2}\Bigr),\qquad R_M=\frac{R_\infty}{1+\mel/M}}
 {$R_\infty=109\,737$~см$^{-1}$; $m$ — номер нижнего уровня серии, $n>m$.}
 {водородоподобные системы.}
 {длины волн спектральных линий и изотопический сдвиг.}

\keyformula{E_{nl}=-\frac{\Ry}{(n-\Delta_l)^2}}
 {$\Delta_l$ — квантовый дефект (зависит от $l$, слабо от $n$).}
 {валентный электрон щелочного атома.}
 {уровни и спектры щелочных металлов; снятие вырождения по $l$.}

\section{Физическая интерпретация}

\textbf{Зависимости от параметров.} $r\propto n^2/(Z\mu)$, $E\propto -Z^2\mu/n^2$: тяжёлая частица
или большой заряд ядра делают систему компактнее и сильнее связанной. Мюонный атом в 207 раз
меньше обычного — он «прощупывает» ядро. Энергии уровней сгущаются к границе ионизации
($E_{n+1}-E_n\approx2\Ry/n^3$), поэтому высоковозбуждённые (ридберговские) атомы имеют огромные
размеры ($r\sim n^2\aB$) и близкие уровни.

\textbf{Связь с экспериментом.} Модель Бора количественно описала серии водорода, объяснила
постоянную Ридберга через фундаментальные константы, предсказала спектр He$^+$ (линии Пикеринга,
ранее приписывавшиеся водороду) и привела к открытию дейтерия по изотопическому сдвигу. Опыт
Франка--Герца (1914) непосредственно подтвердил дискретность энергетических уровней атомов.

\textbf{Ограничения.} Модель — полуклассическая: правило квантования постулировано. Последовательное
описание даёт уравнение Шредингера (\tref{t:10}), которое воспроизводит те же $E_n$, но
вместо орбит даёт распределения вероятности, а в основном состоянии $L=0$.

\section{Что необходимо знать на экзамене}

\begin{examcheck}
\checkhead{Определения}
\begin{itemize}
 \item три постулата Бора; правило квантования $L=n\hbar$ и его связь с волнами де Бройля;
 \item $\aB$, $\alpha$, $\Ry$ — формулы и численные значения; приведённая масса; квантовый дефект.
\end{itemize}
\checkhead{Формулы}
\begin{itemize}
 \item $r_n=\aB n^2/Z$, $v_n=Z\alpha c/n$, $E_n=-Z^2\Ry/n^2$ (с поправкой $\mu/\mel$);
 \item формула Ридберга, серии Лаймана, Бальмера, Пашена;
 \item $\Delta\lambda/\lambda\approx\mel/2M_p$ для H--D; $E_{nl}=-\Ry/(n-\Delta_l)^2$.
\end{itemize}
\checkhead{Выводы}
\begin{itemize}
 \item $r_n$, $v_n$, $E_n$ из закона Ньютона и правила квантования; $E=-T=U/2$;
 \item сведение задачи двух тел к задаче с приведённой массой; изотопический сдвиг;
 \item принцип соответствия при $n\gg1$.
\end{itemize}
\checkhead{Интерпретации}
\begin{itemize}
 \item почему $v\ll c$ (малость $\alpha$) и когда это нарушается ($Z\alpha\sim1$);
 \item физика квантового дефекта (проникновение в остов); масштабирование для ионов и мезоатомов;
 \item почему мезоатомы дают сведения о размерах ядер; недостатки модели Бора.
\end{itemize}
\end{examcheck}
